\part{States, Superposition, and Measurement}

\chapter{Waves force linearity}
\label{ch:waves}

\physicalquestion{Why can two waves occupy the same string at the same time
without destroying one another?}

\modelassumptions{The string is homogeneous, perfectly flexible, held at
constant tension, and displaced transversely by a small amount.  Slopes are
small enough that quadratic geometric corrections, damping, and external
forcing may be neglected.}

\section{What language can express superposition?}

A real \emph{vector space} is a set $V$ with addition and scalar
multiplication satisfying the usual distributive, associative, and inverse
laws.  A map $T:V\to W$ is \emph{linear} when
\[
T(av+bw)=aT(v)+bT(w).
\]
These axioms are not introduced because coordinates are convenient.  They
encode the empirical modeling rule that amplitudes add and may be rescaled.

For a string of length $L$, let $u(x,t)$ be its transverse displacement.  In
the small-slope model, force balance gives
\begin{equation}
\rho\,\partial_t^2u=T\,\partial_x^2u,
\qquad u(0,t)=u(L,t)=0,
\label{eq:wave}
\end{equation}
where $T>0$ is tension and $\rho>0$ is linear density.

\section{What does the wave equation force?}

\begin{theorem}[Superposition]
The solutions of \eqref{eq:wave} form a real vector space.  For every fixed
time $t$, evolution from admissible initial data $(u_0,v_0)$ to
$(u(\cdot,t),\partial_tu(\cdot,t))$ is linear whenever the initial-value
problem has a unique solution.
\end{theorem}

\begin{proof}
Let $D=\rho\partial_t^2-T\partial_x^2$.  Differentiation is linear, so
$D(au+bv)=aD(u)+bD(v)$.  The boundary conditions are homogeneous and are
therefore preserved by the same combination.  Hence solutions are closed
under addition and scaling.  If $E_t$ denotes evolution, then
$aE_t(u_0,v_0)+bE_t(w_0,z_0)$ solves the equation with initial data
$(au_0+bw_0,av_0+bz_0)$.  Uniqueness identifies this solution with
$E_t(au_0+bw_0,av_0+bz_0)$.
\end{proof}

\section{What does this say about an observed string?}

The theorem predicts that independently prepared small-amplitude normal
modes pass through one another and that their measured displacement is the
sum of their displacements.  It does not predict exact superposition for a
large displacement: there the omitted nonlinear terms invalidate the model.
For fixed endpoints the functions
\[
u_n(x,t)=\sin\!\left(\frac{n\pi x}{L}\right)
\cos\!\left(\frac{n\pi}{L}\sqrt{T/\rho}\,t\right)
\]
are modes, and linearity makes every finite linear combination another
motion.  The coefficients are coordinates used to calculate a motion; the
solution space and its linear evolution exist before a basis is chosen.

\section{What further contact should be proved?}

\begin{exercise}[Energy distinguishes the physical solution space]
For sufficiently regular solutions of \eqref{eq:wave}, prove conservation of
\[
E(t)=\frac12\int_0^L\bigl(\rho|\partial_tu|^2
+T|\partial_xu|^2\bigr)\,dx.
\]
Explain which boundary term uses the fixed-end assumption.
\end{exercise}

\begin{exercise}[Critical damping breaks conservative mode evolution]
Write $m\ddot q+b\dot q+kq=0$ as a first-order linear system.  Classify its
three damping regimes by the characteristic polynomial and identify the
nontrivial Jordan block at $b^2=4mk$.
\end{exercise}


\chapter{Composite systems require tensor products}
\label{ch:tensors}

\physicalquestion{If two systems are prepared independently, what state
space contains both independent preparations and correlations between them?}

\modelassumptions{Only the linear structure of each component is retained.
The preparation rule is linear in either component when the other component
is fixed.}

\section{Why are pairs of vectors not enough?}

The direct sum $V\oplus W$ describes alternatives or independent sectors: a
vector is a pair $(v,w)$.  A quotient $V/N$ declares vectors differing by an
element of $N$ equivalent.  Neither construction represents a preparation
that is bilinear in two inputs.

The \emph{tensor product} is a vector space $V\otimes W$ with a bilinear map
$\tau:V\times W\to V\otimes W$ such that every bilinear
$b:V\times W\to X$ factors through a unique linear map
$\widetilde b:V\otimes W\to X$:
\[
b=\widetilde b\circ\tau.
\]
This universal property, rather than a choice of product basis, defines the
composite state space.

\section{What makes the tensor product canonical?}

\begin{theorem}[Existence and uniqueness of the tensor product]
Every pair of vector spaces has a tensor product.  Any two tensor products
satisfying the universal property are related by a unique isomorphism that
preserves pure tensors.
\end{theorem}

\begin{proof}
Take the free vector space on $V\times W$ and quotient by the span of the four
families enforcing additivity and homogeneity in each variable.  The class of
$(v,w)$ is denoted $v\otimes w$.  A bilinear map vanishes on those relations,
so its linear extension descends uniquely to the quotient.  If $(T,\tau)$ and
$(T',\tau')$ both satisfy the property, universality gives maps
$f:T\to T'$ and $g:T'\to T$ preserving pure tensors.  Uniqueness applied to
the identity factorization gives $gf=\id_T$ and $fg=\id_{T'}$.
\end{proof}

\section{Where does entanglement enter?}

A nonzero vector $\psi\in V\otimes W$ is a \emph{product vector} if
$\psi=v\otimes w$; otherwise it is entangled.  The universal property
permits all bilinear product preparations, but the resulting vector space
necessarily contains their linear combinations.  Thus entanglement is not
an added force or interaction.  It is already present in the minimal linear
space closed under superposition of composite preparations.

\section{What further contact should be proved?}

\begin{exercise}[Schmidt decomposition]
For finite-dimensional complex inner-product spaces $H_A,H_B$, associate to
$\psi\in H_A\otimes H_B$ the map $H_A^*\to H_B$.  Apply singular-value
decomposition to prove
$\psi=\sum_i s_i a_i\otimes b_i$ with orthonormal families and $s_i>0$.
Prove that $\psi$ is a product vector exactly when one Schmidt coefficient is
nonzero.
\end{exercise}

\begin{exercise}[Universal constructions]
Give coordinate-free universal properties for direct sums and quotients.
Use them to prove $k(X\sqcup Y)\cong kX\oplus kY$ and
$k(X\times Y)\cong kX\otimes kY$ for free vector spaces.
\end{exercise}


\chapter{Measurement requires Hilbert-space geometry}
\label{ch:measurement}

\physicalquestion{How can a state assign probabilities to mutually exclusive
measurement outcomes without choosing coordinates?}

\modelassumptions{The system is finite dimensional, pure states are rays in a
complex vector space, and an ideal repeatable measurement is represented by a
self-adjoint operator.}

\section{What structure turns amplitudes into probabilities?}

A Hermitian inner product on a complex vector space $H$ is positive definite,
linear in its second argument, and conjugate linear in its first.  It gives
$\|\psi\|^2=\langle\psi,\psi\rangle$.  The adjoint of $A:H\to H$ is the
unique $A^*$ satisfying
$\langle A\psi,\phi\rangle=\langle\psi,A^*\phi\rangle$.
An operator is self-adjoint when $A=A^*$ and normal when $A^*A=AA^*$.

An orthogonal projection is an operator $P=P^*=P^2$.  If $\|\psi\|=1$, the
number $\langle\psi,P\psi\rangle=\|P\psi\|^2$ lies in $[0,1]$ and is
independent of the phase chosen to represent the ray of $\psi$.

\section{Why do observables have spectral outcomes?}

\begin{theorem}[Finite-dimensional spectral theorem]
Every normal operator on a finite-dimensional complex inner-product space has
a unique decomposition
\[
A=\sum_{\lambda\in\sigma(A)}\lambda P_\lambda,
\qquad P_\lambda P_\mu=0\ (\lambda\ne\mu),
\qquad \sum_\lambda P_\lambda=1,
\]
where $P_\lambda$ projects onto the $\lambda$-eigenspace.
\end{theorem}

\begin{proof}
The characteristic polynomial gives an eigenvector $v$.  Normality implies
that $v^\perp$ is invariant under both $A$ and $A^*$: for $w\perp v$, use
$\ker(A-\lambda)^\perp=\im(A^*-\overline\lambda)$.  Induction on dimension
produces an orthonormal eigenbasis.  Grouping basis vectors by eigenvalue
gives the projections.  Their images are the intrinsic eigenspaces, so the
decomposition is unique.
\end{proof}

\section{How does the theorem return to measurement?}

For a self-adjoint observable $A$, all eigenvalues are real.  The ideal
measurement model assigns
\[
\Pr_\psi(A=\lambda)=\langle\psi,P_\lambda\psi\rangle,
\qquad
\sum_\lambda\Pr_\psi(A=\lambda)=1.
\]
The theorem supplies mutually exclusive outcomes and the inner product turns
their components into normalized probabilities.  This is a mathematical
model of measurement statistics; agreement of repeated experiments with
these frequencies tests the model rather than proving the spectral theorem.

\section{What further contact should be proved?}

\begin{exercise}[Uncertainty from noncommutativity]
For self-adjoint $A,B$ and unit $\psi$, set
$\Delta_\psi A=\|(A-\langle A\rangle_\psi)\psi\|$.  Use Cauchy--Schwarz to
prove
\[
\Delta_\psi A\,\Delta_\psi B
\geq\frac12|\langle\psi,[A,B]\psi\rangle|.
\]
State precisely when equality holds.
\end{exercise}

\begin{exercise}[Normal modes are spectral measurements]
For a finite chain of equal masses and springs, express its stiffness as a
positive self-adjoint operator.  Find its spectral projections and recover
the mode frequencies without treating the chosen mass coordinates as
intrinsic.
\end{exercise}


\chapter{Ancillas force complete positivity}
\label{ch:channels}

\physicalquestion{Why must a physical evolution remain positive after the
system is coupled to an arbitrary auxiliary system on which nothing is done?}

\modelassumptions{States of an $n$-level system are density operators in
$M_n(\C)$.  Dynamics is linear on statistical mixtures, and an unused
ancilla evolves by the identity map.}

\section{Why is positivity alone insufficient?}

A density operator is $\rho\geq0$ with $\tr\rho=1$.  A linear map
$\Phi:M_n\to M_m$ is positive when it carries positive matrices to positive
matrices.  It is \emph{completely positive} when
$\id_{M_r}\otimes\Phi$ is positive for every $r$.  A quantum channel is a
completely positive trace-preserving map.

The transpose is positive but $\id_{M_2}\otimes t$ sends the maximally
entangled two-qubit state to an operator with a negative eigenvalue.  Ordinary
positivity therefore fails the untouched-ancilla test.

For $\Phi:M_n\to M_m$, define the Choi matrix
\[
C_\Phi=\sum_{i,j}E_{ij}\otimes\Phi(E_{ij}).
\]

\section{What structure characterizes physical finite-dimensional maps?}

\begin{theorem}[Choi--Kraus--Stinespring]
For a linear map $\Phi:M_n\to M_m$, the following are equivalent:
\begin{enumerate}
\item $\Phi$ is completely positive;
\item $C_\Phi\geq0$;
\item $\Phi(x)=\sum_{a=1}^rK_axK_a^*$ for some $K_a:\C^n\to\C^m$.
\end{enumerate}
Consequently $\Phi(x)=V^*(x\otimes1_E)V$ for a map
$V:\C^m\to\C^n\otimes E$.  If $\Phi$ is unital, $V$ is an isometry.
\cite{choi,stinespring}
\end{theorem}

\begin{proof}
Complete positivity makes
$C_\Phi=(\id\otimes\Phi)(|\Omega\rangle\langle\Omega|)$ positive, with
$\Omega=\sum_i e_i\otimes e_i$.  A positive decomposition
$C_\Phi=\sum_a|v_a\rangle\langle v_a|$ identifies each $v_a$ with a matrix
$K_a$ and direct substitution gives the Kraus formula.  That formula is
positive after every matrix amplification.  Finally set
$V\xi=\sum_aK_a^*\xi\otimes e_a$; then
$V^*(x\otimes1)V=\sum_aK_axK_a^*$.  The identity
$V^*V=\Phi(1)$ proves the last assertion.
\end{proof}

\section{What does dilation say about open evolution?}

In the Schr\"odinger picture the adjoint form of the theorem says that every
channel can be written
\[
\rho\longmapsto\operatorname{Tr}_E(U(\rho\otimes\sigma)U^*)
\]
after enlarging the environment if necessary.  Irreversibility is therefore
compatible with reversible evolution of a larger closed system: information
is lost only when the environment is ignored.  Complete positivity is the
coordinate-free consistency condition that makes this statement survive
every added ancilla.

\section{What further contact should be proved?}

\begin{exercise}[No universal cloner]
Suppose a channel maps $\rho\otimes\sigma$ to $\rho\otimes\rho$ for every
pure state $\rho$.  Use monotonicity of fidelity, or contractivity of trace
distance, to derive a contradiction from two nonorthogonal states.
\end{exercise}

\begin{exercise}[Part I synthesis]
Begin with two normal modes of the string in Chapter~\ref{ch:waves}, regard
their complex span as a two-level system, attach a two-level environment, and
construct a unitary coupling whose reduced evolution is a dephasing channel.
Find its Kraus operators and Choi matrix, prove complete positivity, and
identify exactly which superpositions lose phase coherence.
\end{exercise}

\parttransition{Linear spaces, tensor products, and channels describe states
and evolutions.  They do not identify which transformations preserve the
physical structure or how those transformations label observable sectors.
That missing classification is symmetry.}
